\chapter{总结与展望}

\label{ch:concl}

\section{工作总结}

本文就 Bender 等人\cite{bender2022otsh}在《On the Optimal Time/Space Tradeoff for Hash Tables》中提出的主要结构展开论述，并精简设计出一个哈希系统，该系统把 HashTable 设作入口，内部按照 Facility，Cubby，Slot 这三层来组织元素；设计一个可逆的哈希函数；利用 Router 来维持局部键与槽位之间的映射关系；在槽位处存储经过商压缩之后的 payload；凭借 MiniArray 来保留用于恢复 \(\pi(x)\) 的 MetaEntry；在执行扩缩容操作的时候则采取 active/old 双表逐步过渡的方式。

功能检测覆盖初始化、插入、查询、删除、删后查询与重复插入等路径。性能检测在 \(n\ge 10^4\) 的多组 \((n,N,K)\) 配置下开展：命中查询 p50 在 $2.9$--$5\,\mu$s 量级，插入 p50 约 $10$--$15\,\mu$s ，insert\_moved\_max 与 router\_steps\_max 均不超过 6，与第三章有界路由及工程版 k-kick 的设计一致。

\section{不足与局限}

本文存在一些不足与局限之处，与论文中的完整理论结构相比仍有所差距，主要体现在以下方面：
\begin{enumerate}
  \item k-kick 机制：当前实现已在单个 Cubby 内达成有界踢出链，但尚未完整模拟论文中的 k-kick 树层次结构；实验中的 kick\_count 反映的是代码层面的搬移次数，不能直接等同于理论交换成本。
  \item Facility 层重建：虽具备多 tiered Cubby、tail 维护及拆分/合并接口，但调度器每次执行完整重建，未将 \(O(|I|)\) 级任务细分为严格的后台步骤；tiered 上限与容量参数仍属工程设定，尚不足以支撑理论均摊界的论证。
  \item MiniArray 与 Router 更新：两者保留了变长编码与局部路由的接口语义，但更新分别依赖全局重打包与整树重建，未达到论文所要求的常数时间局部更新。
  \item 空间统计：meta\_bits\_total、Router 编码位数等指标仅反映逻辑编码规模，未计入 vector 指针、对象头等运行时内存开销，不能作为严格空间上界的证明依据。
\end{enumerate}

实验范围仍存在局限：虽已在 \(n\in[1.5\times10^4,10^5]\) 上扫过多组 \((N,K)\)，但未变化负载因子与随机种子分布，也未与标准库、IcebergHT\cite{pandey2023iceberght} 或其它紧凑哈希实现\cite{kurpicz2023pachash,lehmann2024shockhash} 对比，结论仅适用于当前实现策略下的行为刻画，不能外推为通用哈希表的最优性能。

\section{可改进方向}

针对上述不足，后续工作可从以下方向推进：
\begin{enumerate}
  \item 在 Cubby 内完整实现 k-kick 树结构，使其与论文中的层级控制及失败处理机制相衔接。
  \item 细化 Facility 层后台重建调度，将拆分/合并拆为可分步执行的任务，并重新校准 tiered 上限、\(t_j\) 与 Cubby 容量公式。
  \item 改进 MiniArray 与 Router 的增量更新策略，降低重打包与整树重建带来的开销，并区分逻辑编码位数与实际内存占用。
  \item 扩展测试覆盖面：增加负载因子、随机种子等控制变量，补充与标准实现的对比实验及可视化分析，更充分地刻画系统的时间与空间行为。
\end{enumerate}